\section{Conclusion and Future Work}
This paper has presented the design, implementation, and empirical evaluation of a lightweight Digital Rights Management (DRM) system optimized for secure video distribution. By leveraging a syntax-aware Selective Encryption (SE) mechanism that targets only critical H.264 structural foundations (IDR, SPS, and PPS units), the system dramatically lowers client-side cryptographic overhead. Experimental profiles confirmed up to a 93.04\% reduction in processing latency, alongside a drop in CPU core utilization to on smaller video profiles, while maintaining effective perceptual confidentiality against unauthorized playback attempts.

Future investigations will expand upon this framework across multiple systemic directions:
\begin{enumerate}
    \item \textbf{Dynamic Codec Parsing Adaptability:} Engineering adaptive extraction loops to natively support H.265/HEVC block layer architectures and royalty-free AV1 temporal grouping parameters.
    \item \textbf{Over-the-Top Protocol Integration:} Encapsulating the selective AES-CTR stream directly inside Dynamic Adaptive Streaming over HTTP (DASH) and HTTP Live Streaming (HLS) operational manifests.
    \item \textbf{Hardware-Isolated Key Architecture:} Moving volatile key initialization states into Trusted Execution Environments (TEEs) using Android Widevine Level 1 or ARM TrustZone equivalents to protect endpoints from local memory inspection vectors.
    \item \textbf{Watermarking Integration:} Layering robust, lightweight visual watermarking markers on top of unencrypted B/P slices to trace illegal distribution pathways back to verified viewer streams.
\end{enumerate}